
\begin{longtable}{p{3cm} p{11cm}}

\textbf{Alpha} & Pita frekuensi alpha: komponen gelombang EEG pada rentang 8--13\,Hz yang dominan pada wilayah oksipital saat kondisi rileks dengan mata tertutup. \\

\textbf{Artefak} & Komponen sinyal non-neural pada rekaman EEG yang berasal dari aktivitas fisiologis (kedipan mata, gerakan otot) maupun gangguan eksternal seperti interferensi perangkat elektronik. \\

\textbf{Augmentasi} & Augmentasi data (\textit{data augmentation}): teknik penambahan variasi data pelatihan secara artifisial untuk meningkatkan performa model dan mengurangi \textit{overfitting}. \\

\textbf{ASD} & Autism Spectrum Disorder: gangguan perkembangan saraf yang ditandai hambatan komunikasi, interaksi sosial, serta perilaku repetitif. \\

\textbf{Beta} & Pita frekuensi beta: komponen EEG pada rentang 13--30\,Hz yang berkaitan dengan aktivitas kognitif aktif, konsentrasi, dan pemecahan masalah. \\

\textbf{Decoder} & Komponen model generatif yang memetakan variabel laten untuk merekonstruksi sinyal EEG sintetis. \\

\textbf{Delta} & Pita frekuensi delta: rentang 0.5--4\,Hz yang dominan pada kondisi tidur nyenyak dan proses pemulihan fisiologis. \\

\textbf{DNN} & Deep Neural Network: jaringan saraf tiruan dalam dengan beberapa hidden layer untuk mempelajari pola kompleks. \\

\textbf{EEG} & Electroencephalography: metode perekaman aktivitas listrik otak secara non-invasif menggunakan elektroda pada kulit kepala. \\

\textbf{ELBO} & Evidence Lower Bound: fungsi objektif pada VAE yang terdiri dari reconstruction loss dan KL divergence. \\

\textbf{Encoder} & Encoder: komponen model yang memetakan data input ke ruang laten berdimensi rendah berbasis distribusi probabilistik. \\

\textbf{FFT} & Fast Fourier Transform: metode transformasi sinyal dari domain waktu ke domain frekuensi. \\

\textbf{FID} & Fréchet Inception Distance: metrik untuk mengukur kesamaan distribusi antara data asli dan data sintetis. \\

\textbf{Gamma} & Pita frekuensi gamma: sinyal EEG dengan frekuensi $>30\,Hz$ yang terkait proses kognitif tingkat tinggi. \\

\textbf{Generalisasi} & Kemampuan model mempertahankan performa pada data yang tidak terlihat selama pelatihan. \\

\textbf{Hyperparameter} & Parameter yang ditentukan sebelum pelatihan seperti learning rate atau jumlah estimators. \\

\textbf{ICA} & Independent Component Analysis: metode pemisahan sumber sinyal untuk mengurangi artefak EEG. \\

\textbf{KL} & Kullback--Leibler Divergence: regularisasi pada VAE untuk menjaga distribusi laten mendekati distribusi normal. \\

\textbf{Latent} & Latent Space: ruang representasi berdimensi rendah yang menangkap struktur penting data. \\

\textbf{mRMR} & Minimum Redundancy Maximum Relevance: metode seleksi fitur berbasis relevansi maksimum dan redundansi minimum. \\

\textbf{Overfitting} & Kondisi ketika model terlalu menyesuaikan data pelatihan sehingga menurunkan performa pada data baru. \\

\textbf{PSD} & Power Spectral Density: distribusi energi sinyal terhadap frekuensi. \\

\textbf{Preprocessing} & Pra-pemrosesan: tahap persiapan data EEG seperti filtering, normalisasi, dan segmentasi. \\

\textbf{SWD} & Sliced Wasserstein Distance: metrik untuk mengukur kesamaan distribusi data melalui proyeksi satu dimensi. \\

\textbf{Theta} & Pita frekuensi theta: rentang 4--8\,Hz yang terkait memori dan kondisi mengantuk. \\

\textbf{U-shaped} & U-shaped spectral profile: pola spektral dengan peningkatan daya pada frekuensi rendah dan tinggi serta penurunan pada frekuensi menengah. \\

\textbf{VAE} & Variational Autoencoder: model generatif berbasis encoder-decoder untuk mempelajari ruang laten probabilistik. \\

\textbf{XGBoost} & Extreme Gradient Boosting: algoritma ensemble tree boosting untuk klasifikasi dan regresi yang efisien. \\

\end{longtable}